% p303_2; source PDF page 303; manually verified editable table.
\begin{tikzpicture}[x=1cm,y=-1cm,inner sep=0pt,outer sep=0pt]
\path[use as bounding box] (0,0) rectangle (16.000,10.745);
\draw[line width=.4pt] (0.000,0.000) rectangle (2.900,0.815);
\node[align=center,text width=2.640cm,font={\linespread{1}\fontsize{10}{13.50}\selectfont\bfseries}] at (1.450,0.407) {误解的类型};
\draw[line width=.4pt] (2.900,0.000) rectangle (7.300,0.815);
\node[align=center,text width=4.140cm,font={\linespread{1}\fontsize{10}{13.50}\selectfont\bfseries}] at (5.100,0.407) {具体情形};
\draw[line width=.4pt] (7.300,0.000) rectangle (16.000,0.815);
\node[align=center,text width=8.440cm,font={\linespread{1}\fontsize{10}{13.50}\selectfont\bfseries}] at (11.650,0.407) {范例};
\draw[line width=.4pt] (0.000,0.815) rectangle (2.900,10.745);
\node[align=left,text width=2.640cm,font={\linespread{1}\fontsize{10}{13.50}\selectfont}] at (1.450,5.780) {误解一：只要客观上具备数额，就定罪量刑};
\draw[line width=.4pt] (2.900,0.815) rectangle (7.300,4.530);
\node[align=left,text width=4.140cm,font={\linespread{1}\fontsize{10}{13.50}\selectfont}] at (5.100,2.672) {（1）只要客观上达到“数额较大”，就定罪};
\draw[line width=.4pt] (7.300,0.815) rectangle (16.000,4.530);
\node[align=left,text width=8.440cm,font={\linespread{1}\fontsize{10}{13.50}\selectfont}] at (11.650,2.672) {农民甲以为自己盗窃的是普通葡萄，实际是高科技试验品，价值40万元。有意见认为，既然该葡萄是数额特别巨大的财物（当然也算数额较大的财物），因此甲构成盗窃罪。然而，由于甲没有认识到该葡萄是数额较大的财物，因此不构成盗窃罪};
\draw[line width=.4pt] (2.900,4.530) rectangle (7.300,10.745);
\node[align=left,text width=4.140cm,font={\linespread{1}\fontsize{10}{13.50}\selectfont}] at (5.100,7.637) {（2）只要客观上达到“数额（特别）巨大”，就适用“数额（特别）巨大”的法定刑};
\draw[line width=.4pt] (7.300,4.530) rectangle (16.000,10.745);
\node[align=left,text width=8.440cm,font={\linespread{1}\fontsize{10}{13.50}\selectfont}] at (11.650,7.637) {客人王某带一盆兰花入住五星级酒店时，服务员甲认为这盆兰花应该值几千元，便盗窃到手。物价部门鉴定，该兰花属于稀世珍品，价值100万元。有意见认为，既然该兰花是数额特别巨大的财物，因此，对甲应适用“数额特别巨大”档次的法定刑。\\{}然而，甲虽然认识到该兰花是数额较大的财物，但是没有认识到是数额特别巨大的财物，因此，只能适用“数额较大”档次的法定刑，不能适用“数额特别巨大”档次的法定刑};
\end{tikzpicture}
